\section{Validation Program}

Item development should begin from behaviorally specific functions rather than evaluative labels. Candidate items should cover attempts to change contact, timing, availability, boundaries, interaction structure, or information flow without presupposing coercion.

The scale-development program requires independent development and confirmation samples, comparison with adjacent constructs, criterion validity against observed or repeatedly measured regulatory behavior, and testing of incremental validity beyond attachment anxiety and avoidance. A failure to establish discriminant validity should lead to revision or abandonment of the construct rather than post hoc redefinition.

Items must specify their reference period and relationship context. Content review should identify redundancy with agency/assertiveness, generalized need for control, intolerance of uncertainty, attachment insecurity, dominance and coercive control before factor analysis. No factor or total score is assumed in advance. Adjacent constructs require authorized instruments and their own measurement-quality assessment; a weak comparator cannot establish distinctness.

Development data support item revision and exploratory factor analysis. An independent confirmation sample tests the frozen structure and scoring rule. Retest evaluates temporal stability; invariance analyses depend on adequate subgroup information and cannot be inferred from pooled fit. Primary endpoints, numerical decision rules, sample-size justification, missing-data procedures and multiplicity remain to be prospectively specified; this text is a research program, not a completed validation protocol.

Criterion validation requires information beyond reworded repetitions of the scale's own items. The initial criterion is a brief independently coded response to standardized neutral vignettes, with opportunity and relationship context recorded. Such responses concern hypothetical intentions and cannot establish enacted daily behavior. Longitudinal ecological measures and private partner reports belong exclusively to M04; no daily diary is included in M03. Incremental prediction must compare relevant adjacent constructs and use an untouched validation sample or prospectively reserved test set.

Planned result sections concern content validity, sample flow, item distributions, exploratory and confirmatory structure, reliability/retest, discriminant validity, invariance where supported, and incremental criterion prediction. They do not include generator-recovery matrices or P/PI/PID parameter-recovery claims. If distinctness fails, report that result and restrict or abandon the proposed construct without redefining it on the final test data.
